\section{Discussion and Limitations}
\label{sec:discussion}

\system{} does not repair a poor global join order.  A locally efficient join
can still materialize an intermediate relation that a different order would
avoid.  The operator can provide observations to a higher-level reoptimizer,
but coordinating that feedback is outside this work.

The confidence model also assumes that observed prefixes are informative after
stratification by storage partition.  Correlated physical order can delay a
heavy key until after commitment.  Conservative upper bounds and the buffer cap
preserve correctness and memory safety, but not optimality.  Randomized page
sampling or storage-level summaries may reduce this risk at additional I/O or
maintenance cost.

Very high output cardinality is another boundary.  Splitting a heavy key
balances the work needed to enumerate matches but cannot avoid the intrinsic
$f_R(x)f_S(x)$ output.  The scheduler must reserve downstream capacity or spill
output.  Finally, the metadata service is logically centralized.  Sharding
improves throughput, but a production system must quantify failover pauses and
garbage collection across many concurrent queries.
